\chapter{Limitations of the Study}



As with most empirical studies, the scope of this work is defined by several
methodological boundaries that should be considered when interpreting the
results. These boundaries do not reduce the relevance of the findings, but
rather define the context in which they should be understood.

First, the experiments focus on Java-to-C++ translation. Although this
language pair is representative because of the syntactic and semantic
differences between the two languages, the findings may not generalize to other
language pairs. Previous studies suggest that results may not transfer equally across all
language pairs, particularly when programming paradigms differ
\cite{llm_limitations}. Future work should therefore evaluate additional
programming language pairs to better understand how model performance varies
across different paradigms, typing systems, and runtime environments. Expanding
the evaluation to multiple language directions would provide a more
comprehensive picture of the generalization capabilities of code translation
models.

Second, the study evaluates two selected code-specific models: StarCoder2 and Qwen2.5-Coder. Although
these are strong open-source code models, other models such as CodeLlama,
DeepSeek-Coder, or larger proprietary systems may exhibit different behavior.
Prior work suggests that larger models generally achieve better functional
correctness, but they also require substantially more computational resources
\cite{llm_survey}. This means that the observed results may partly reflect the
characteristics of the selected model size rather than the overall capability of
large language models. Future research should include a broader range of models,
particularly larger parameter models and alternative architectures, in order to
better understand the relationship between model scale, architecture, and
translation quality. Such comparisons would help determine whether the observed
performance trends are model-specific or general across different families of
code models.

Another methodological boundary is the use of XLCoST as the training dataset. While XLCoST provides aligned Java and C++ code examples, it
still captures a broad but not exhaustive range of coding styles, programming
domains, and problem types. Dataset quality and diversity are known to have a strong impact
on model performance, especially for code-related tasks \cite{zhu2022xlcost}.
Future work may benefit from using larger and more diverse datasets, including
code from different domains, programming styles, and levels of complexity.
Incorporating additional data sources or combining multiple datasets could
improve the robustness and generalization of fine-tuned models.

The evaluation setup is also defined by the use of CodeEditorBench as the
execution-based benchmark. Although
CodeEditorBench provides realistic evaluation scenarios, other benchmarks such
as xCodeEval \cite{xcodeeval}, CodeTransOcean \cite{codetransocean},
RepoTransBench \cite{repotransbench}, TRACY \cite{tracy_benchmark} or
PolyHumanEval \cite{mhumaneval} may highlight complementary aspects of model behavior, as they emphasize different tasks and evaluation criteria. Therefore, the conclusions
of this thesis are limited to the specific benchmark setup that was used. Future
research should incorporate multiple benchmarks and evaluation methods in order
to obtain a more comprehensive assessment of translation quality. Combining
similarity-based metrics with multiple execution-based benchmarks provides a more robust and reliable evaluation.

Finally, the experiments are conducted under practical hardware constraints. For this reason, the study focuses on 3B-parameter models and uses
parameter-efficient fine-tuning techniques. Larger
models and longer training schedules may have achieved better results, but they
would require significantly more computational resources \cite{hu2021lora}.
Future work should explore the use of larger models, extended training
schedules, and more computationally intensive fine-tuning strategies. Advances
in hardware and optimization techniques may enable more effective training of
large-scale code models in future studies.